\SourceLeaf{414}%
From \(-\frac12b[2]+\frac12a[3]\) and \(\frac12B[2]-\frac12A[3]\), if we write for brevity
\[\begin{gathered}aB-bA=\lambda,\quad ab(A+B)=f,\quad\tfrac12AB(A+B)+\lambda n=g\\
AB(a+b)=F,\quad\tfrac12aBB+\tfrac12bAA=G\end{gathered}\]
we obtain
\[\begin{gathered}\lambda T+G=\lambda(t+f)x,\quad\lambda t-g=(\lambda T-F)x\quad\text{or}\quad\lambda(t-Tx)=g-Fx\end{gathered}\]
and consequently
\[\begin{aligned}\lambda t&=\frac{+g-(F+G)x+fxx}{1-xx}\quad.\ .\ .\ .\ .\ .\ .\ .\ [4]\\
\lambda T&=\frac{-G+(f+g)x-Fxx}{1-xx}\quad.\ .\ .\ .\ .\ .\ .\ .\ [5]\end{aligned}\]
\par
Equation \([1]\) in the form \((t-Tx)^2+TT(1-xx)=rr\) gives
\[\lambda\lambda rr-(g-Fx)^2=\lambda\lambda TT(1-xx)\quad.\ .\ .\ .\ .\ .\ [6]\]
\par
Substituting the value of \(\lambda T\) from \([5]\) into this equation gives the cubic equation
\[\begin{aligned}&2fFx^3-(ff+2fg+FF+2FG+\lambda\lambda rr)xx\\
&\quad+(2fG+2gF+2gG)x+\lambda\lambda rr-gg-GG=0\quad.\ .\ [7]\end{aligned}\]
\par
Once \(x\) has been determined in this way, we find the coordinates of point \((2)\) from the expressions above. If we substitute into them the values just given for \(t-Tx\) and \(T\), they take the following form
\[n-\frac{g-Fx}{\lambda},\qquad\frac{-G+(f+g)x-Fxx}{\lambda\BareRoot{(1-xx)}}\]
and we find the other two points \((1)\) and \((3)\) by adding \(-Ax,-A\BareRoot{(1-xx)}\) and \(+Bx,+B\BareRoot{(1-xx)}\), respectively, to these values.

Since two values of the angle correspond to each cosine, or equivalently since the radical \(\BareRoot{(1-xx)}\) may be taken positive or negative, each admissible root of the equation gives two solutions, namely two positions, symmetric about line \(\mathrm{(II)}\), of the points \((1),(2),(3)\), as is already clear in itself. If \(+1\) or \(-1\) is a root of equation \([7]\), however, the formula above for the ordinate cannot be used, because numerator and denominator then both vanish, and instead we must
